\section{难点，以及我是怎么改的}

这一节按“我本来以为是什么样 → 哪里不对 → 改了什么 → 现在是什么数”来写。选的是真卡住过我的地方。

\subsection{测试全绿，功能没生效}

这是我最想讲的一条，因为同一个形状在一天里出现了六次：代码看着对、测试全绿、功能根本没生效。

\begin{itemize}[leftmargin=1.6em,itemsep=0.25em,topsep=0.3em]
\item 服务端校验上下文模式的白名单是 \code{camp / pve / pvp-live}，加本地对战时客户端开始发 \code{pvp-local}，白名单没跟着改。本地对战里的模型请求全部被 400 拒绝、静默回退成本地规则。规则建议照常工作，界面上没有任何异常——这个模式从头到尾没有产生过一次模型解释。
\item 新增的 \code{rules.js} 没进服务端的静态资源白名单，\code{/rules.js} 返回 404，\code{app.js} 的 import 图整体失败，\textbf{整页白屏}。
\item 陪练的 \code{notify()} 定义了，但全仓库没有一个调用点。主动气泡写完了，从不出现。
\item 军师改造留下一个未定义变量，每秒抛一次 \code{ReferenceError}，浏览器日志里攒了 75 条。
\item 对方面板把同一段内容重复显示了三处。
\item 敌方补位“决定了却无人提交”：AI 选好了合法换宠、存进变量就返回，而玩家那条路径又拒绝替它提交。结果是界面永久冻结在“正在出招…”。
\end{itemize}

还有一个同类的：犹豫判定用的是墙钟时间，玩家中途去操作对方面板的 13 秒被算成了“犹豫了 13 秒”。

单元测试查不出这些，因为它们测的是函数本身，不测函数有没有被接上。所以补了两层检查。

第一层是 \path{wiring.test.js}，离线、确定性，查三件事：\code{app.js} 里每个 \code{\$('id')} 都必须在 \code{index.html} 或 \code{app.js} 自己建出来的元素里存在；\code{app.js} 的死函数清单必须恰好是两个已知遗留，多一个就失败；面试题要求的每项能力（军师局内提示、整局复盘、小测、聊天面板、陪练主动气泡）都必须在 \code{app.js} 里找得到入口。它当场就找出一处指向已删除元素的残留引用。

第二层是 \code{npm run test:smoke}，headless Chrome 真跑一局训练赛，断言营地渲染出 12 张伙伴卡（说明模块图活着）、比赛能打到结果、控制台零报错。变量未定义这类错误只有真页面看得见。它故意不进 \code{npm test}——需要浏览器的检查不该当成所有人的必过门。

这一层也有它抓不到的东西，我自己就见过一次假阳性。冒烟测试第一次跑出“最长连续无可点 14 次”，我以为快卡死了。查下去发现是它从没设置播放速度，每一回合都在放动画、按钮本来就禁用——这个指标量的是动画时长，不是死锁。改成先设最快速度、并且只在横幅不显示“正在”的时候计数，同一个场景现在是 0。

\subsection{玩家不打开聊天，就等于没有教练}

第一版做出来我自己试了两局，发现一件事：小芽的所有能力都在聊天面板里，而我全程没打开过那个面板。一个要玩家主动召唤才存在的助手，在真实游玩里等于不存在。

改法是把主入口从聊天框挪到游戏事件上：开局和关键风险变化时，战斗页底部出现一行短提示，解释收在二级入口里按需展开；玩家出招之后，旧提示立刻隐藏，等新局面重新评估。聊天框降级成补充入口。

副作用是提示一多就变成噪音。所以门控最后比内容更难做：玩家是否在前台、是否已经出招、是否已经忽略过、这个回合是否已经提示过、上次提示距今多久、当前风险是否构成可行动的风险，逐项判断；每局总量、同类提示只出一次、冷却时间各有上限，老师和军师还共用同一份每局配额。

这里有一条我从教育数据挖掘的文献里学来的规则。那篇讲 assistance dilemma 的论文里，作者报告他们早期的预测器没有把“学生用了提示”这件事算进去，结果因为提示本身会把人推向高效步骤，用过的学生反而被判定成“不需要帮助”，他们的修正是把两者分开。

小芽里是同一件事：判断玩家是否掌握某个知识点时，只用他没有被提示过的时候做出的选择，代码里就是一个过滤条件。我是先按直觉这么设计的，后来才发现这个坑在文献里有记录。

\subsection{数字不能让模型算}

我一开始让模型看局面然后给建议，试出来的回答里数字是错的。不是偶尔错，是稳定地错：编出不存在的技能名，算错伤害，在两个都能击倒的技能里挑面板数字大的那个。

后来把职责切开——引擎枚举、模型解释，也就是前面军师那一节讲的架构。这之后还剩两类错误是数字校验抓不到的。

第一类是道具名。游戏里只有回复药、净化药、能量果三样，模型会把净化药说成“解药”、能量果说成“以太”。数字全对，引用全对，只有名字是错的，因为没有数字可校验。

第二类是因果。宠物倒下时，它这一回合已经排好的行动会被取消，模型会把这次取消描述成“命中了”。同样没有数字可比。

这两类都靠加检查解决：维护一份别名黑名单，命中就判不合格；读回合事件，某个行动被记为取消时，正文里就不能出现这个行动造成了伤害。判不合格之后走既有的降级路径，显示规则结论而不是错误描述。

顺带说一处同类的文案问题。降级提示原来会把服务端原文（“教练上下文无效”）和内部错误码直接端到玩家面前。现在按原因映射成四句玩家话，内部代码只进日志。后台术语那件事我人工扫了三轮，三轮都漏：第一轮漏了若干处，第二轮又漏 5 处，第三轮才发现知识卡自己的 principle 正文里也有一个。靠一次扫除是扫不干净的。

所以这轮把它改成了机械检查。\code{evals/player-copy.test.js} 遍历玩家能读到的每一个字符串——军师短句、点开“看看原因”后的依据、被引用的知识卡正文、阵容建议、局内提示、收尾提示、整局复盘——命中禁止词表（“启发式”“枚举”“分支”“平均收益”这类）就失败。我做了负向验证：把“启发式”植入一句可达文案，测试失败；移除，测试通过。这一步是必须做的，因为\textbf{一个不会失败的检查比没有检查更糟}——它会让人以为这件事已经守住了。

这个检查现在还有一个缺口，我写清楚以免它被当成完备的：它用 \code{observe()} 取局内提示时没有传 incident，所以“明显策略错误”那条理由对应的正文不在覆盖范围内，而那段正文里仍然留着“一回合启发式评分”。人工扫的第三轮也漏在这里，同一个位置。修它需要把 incident 造进测试，不在这轮的范围里。

\subsection{“该不该调工具”不该由模型决定}

工具决策我做过三轮真实评测，同一个预注册用例集、44 条、全部真实 API 调用，每轮都保留原始数据，包括变差的那一轮。三轮依次改的是：原始提示词（“先检查已有证据，再决定是否补证据”）；把“默认是停止”写进提示词并点名哪几类必须调；以及给模型加一个“必须调用”的字段。第 3 轮已回退，第 2 轮是仓库当前的状态。

\begin{table}[h]
\centering
\small
\begin{tabular}{L{3.6cm}L{2.9cm}L{3.4cm}L{3.4cm}}
\toprule
指标 & 第 1 轮 & 第 2 轮 & 第 3 轮 \\
\midrule
多余工具调用率 & 75.0\% & 33.3\% & 13.6\% \\
漏调率 & 0.0\% & 21.9\% & 77.3\% \\
严格正确率 & 31.25\% & 50.0\% & 31.8\% \\
模型自停率 & 3.23\% & 96.8\% & --- \\
延迟 p50 & 4284 ms & 2798 ms & 2154 ms \\
44 条成本 & \$0.0335 & \$0.0209 & \$0.0178 \\
\bottomrule
\end{tabular}
\end{table}

原文里“先检查再决定”这句话把“去查”当成了默认动作——12 条本该零调用的用例里 9 条调了工具，31 次决策里只有 1 次是模型自己判断停下来的。

第 3 轮为什么变差值得单独说。我加那个字段本意是补上三类硬需求，实际上等于告诉模型“这个字段是 null 就不必调”——而它在大多数用例上是 null。整个判断被推给了那个字段，漏调因此从 21.9\% 跳到 77.3\%。这是典型的过度规定：用一个新字段表达“必须调”，同时无意中表达了“其余都不必调”。

三轮的摆动说明问题不在措辞，在分工。数据给出的分界很清楚：模型\textbf{选哪个工具}很准（第 1 轮 90.63\%；第 2 轮在“预期应调且确实调了”的 12 条上是 12/12），但\textbf{判断该不该调}只有一半正确。既然一半是错的，就不该让模型做这个判断。所以第 4 轮不是改提示词，是把政策写进代码（\path{coach/runtime.js} 里的 \code{policyFor}）：证据包里已经带着当前局面和最近回合事件，问当前状态的不用再查；只有结构上不含的四类才需要查。判定需要时，代码直接执行第一步，之后再问模型要不要继续查；判定不需要时，整个工具循环跳过。模型现在只在“还要不要继续看”这个方向上参与。

边界要说清。44 条里只有 32 条可判定，每格 1 条、单次运行、没有重复测量，所以上表所有差值都没有置信区间；第 2 轮和第 3 轮之间 18 个百分点的差距，有多少是真实效应、多少是单次波动，这批数据回答不了。而且严格正确率即使在第 2 轮也只有 50\%。另外还有一处规格分歧我主动交代：评测的预注册标准认为“我现在场上这只还剩多少血”应该调一次 \code{read\_state}，但证据包里本来就带着当前局面，调它只是把已有的事实再取一遍，第 2 轮那 7 条“漏调”里有 6 条属于这一类。所以“50\%”里有相当一部分不是模型错，是两种政策的差异。产品必须选一个并说清楚，我选了“包里已有的事实不重复取”；按新口径重新标注必须在看到结果之前完成，这一轮我没做。

\subsection{检索没有赢过关键词}

我做过一个假设：把检索从关键词换成“关键词加多语言向量”的混合方案，口语化提问的命中率会提高。为了验证它，我建了一个四臂对照：不检索、纯关键词、混合，外加一个把反例文本提为一等字段的第四臂。128 条自查查询，112 条有正解、16 条越界，离线跑，不调用模型。

结果是没有一个臂显著赢过关键词基线。Hit@3 分别是 0.00\%、88.39\%、90.18\%、91.07\%，最好的第四臂相对关键词只高 2.68 个百分点，McNemar 精确检验 $p = 0.4531$，配对 bootstrap 95\% 置信区间 $[-1.79, +7.14]$ 个百分点，跨零。按这个不一致对的比例，要有 80\% 把握检出这个量级的差异，大约需要上千条查询。

这里我差点犯一个错。看到混合没赢，我的第一反应是“检索没用，RAG 是伪需求”。把四臂拆开看才发现不是：第四臂的 MRR 从 0.8007 升到 0.8244，排序确实好了。但这些改善都发生在 112 条正解上，一条查询的变动就能让 Hit@3 动将近一个百分点。所以我没有拿它证明任何事——既没有换成向量检索，也没有宣称 RAG 无效，而是记成“这次测量没测出差别”，把关键词保留为默认。

真正被这轮测出来的问题是负例拒答：三个检索臂都只有 56.25\%（9/16）。“今天晚饭吃什么”“1+1 等于几”这类明显越界的提问仍然召回了卡片，原因是词项检索只要命中任意二元组就判为正分，没有最低分阈值，也没有拒答机制。这是未解决的问题，而且调权重解决不了。

还有一条边界必须写在前面：这 128 条查询是我自己写的，标注也是我自己读卡片文本后打的，不是独立盲标集，没有第二标注者，也没有一致性检验。一个直接证据是，沿用下来的那 20 条旧查询 Hit@3 是 81.25\%，这次新写的 92 条是 93.48\%——新写的明显更高，说明我出题时已经知道卡片里有哪些词。所以上面所有绝对数字都偏乐观，不能当外部基准。

\subsection{模型会说过期的话}

表现是：玩家已经出招了，模型那边的解释才回来，界面上出现一句针对已经过去的局面的建议。

我试过在提示词里写“如果局面已经变化就不要回答”，没用。模型看不到自己那句回答会在什么时候被显示出来，这个判断只能由程序做。

改法是把每个请求绑定到具体的局面版本：请求发出时记下对局 ID、回合号、规则版本和上下文 epoch，返回时逐项比对，任何一项变了就丢弃结果，不显示也不播报。客户端断连时服务端同时中止上游请求。这条有 6 项可复现的测试，用替换请求的方式构造慢响应，验证“快速出招后不显示旧文字、不重复补发”。真实浏览器里叠加慢响应这一半没做，需要可控的慢网环境，和语音的取消验证一起挂着。

\subsection{对手原来不交给模型，后来交给了模型}

题目场景是 PVP，我一开始想让模型来当对手，理由是“这样显得 AI 更强”。查了资料之后放弃了这个想法。

放弃的理由不是我懒，是证据方向相反。PokéLLMon 是一个纯靠语言模型打宝可梦的工作，阶梯赛胜率 49\%。PokéChamp 换成“LLM 辅助 minimax 搜索”之后，对最强规则型 bot 胜率 84\%，而且不需要额外训练。Metamon 进一步用离线强化学习加 Transformer，论文里说它的表现超过了 LLM 智能体和强启发式搜索引擎。在这个游戏类型里，决定胜负的是搜索和训练，不是语言能力。所以对手出招我用了引擎的启发式：轻松档随机、标准档优先伤害和治疗、挑战档枚举双方合法行动并比较一回合结果。

这一轮我把它改了：\textbf{对手的决策交给模型，引擎只负责列出有哪些合法行动}。上面那段调研的结论没有变，变的是它回答的问题——它说的是“谁更强”，而这轮要解决的是“对面得是个人”。一个按伤害公式调出来的启发式，出招永远挑不出毛病，也永远不像人；难度只是它的参数，不是它的性格。所以引擎留下来做两件事：用 \code{legalActions} 给出合法行动列表，以及在不该由模型决定的时候兜底。真正做选择的换成了模型，从那份列表里挑一个。

\code{chooseEnemy()} 因此降级成纯兜底，只在四条路径上出现：模型未配置（约 1 毫秒返回）、模型选了列表外的行动、超时（4 秒封顶）、网络错。这四条各自有测试，而且都必须能在没有后端、没有密钥的前提下工作。

难度改成只控制“给对手多少帮助”，不控制谁在决策。简单档只给合法行动列表；标准档多给一句对手自己教练的建议；挑战档再给出这一回合的枚举分数。同一个模型，三档拿到的是三份信息量不同的简报。这里有个复用值得一提：对手那个“自己的教练”不是另写一套，是拿同一份军师推理函数、把 player 和 enemy 对调后跑出来的，所以它读到的是自己的血量、能量和技能。对手和玩家用的是同一个大脑，只是坐在对面。

对手的教练在界面上仍然显示，但做了\textbf{高斯模糊}：看得见它存在，读不出内容，像隔着毛玻璃看别人的屏幕。真人同机分屏时不糊——那本来就是两个人共用一块屏。这里要说清一件事：模糊是观感手段，不是安全边界。对手教练的建议本来就基于双方公开的行动，不是隐藏信息，把它糊掉只是不想让玩家读着一句话打牌，不是防作弊。

\textbf{实测}（真实模型，服务在 8765，四次对照合计 543 次决策）：往返 p50 约 1.05 秒、p90 约 1.25 秒、最大 1.6 秒。玩家按正常节奏思考时\textbf{每回合等待 0 毫秒}——对手思考的那一秒被玩家自己的思考时间吸收了，界面上不出现“对手正在思考…”。玩家秒点时才会真等到，实测 0.89 到 2.52 秒。

兜底路径里只有“上游超时”真的触发过，而且只在并发压力下触发：另一批批跑时 135 次请求有 36 次撞上 4 秒上限，兜底率从 0 跳到 27\%，延迟 p90 被抬到 3.7 秒。所以这条链路的真实上限不是“模型有多快”，而是“上游在并发下有多稳”，4 秒超时是它唯一的安全垫。模型未配置、选择非法、网络错这三条一次都没被真实流量碰到，目前只有单测保证。

\textbf{胜率}这一项我没测出结论，原因比“样本小”更具体。挑战档下 agent 与引擎的选择 93\% 相同（142 次里 10 次不同），标准档下 48\% 不同（142 次里 74 次不同，两轮独立复跑都是这个量级）；但把 24 局放在一起看，\textbf{对手一局都没赢过}——四个格子两边的胜率是 100/100、100/100、66.7/50、50/50，这个阵容下胜负根本分不开两臂。所以准确的说法是：n=6 而且对手全败，这套对照既不能支持“对手变强了”，也不能支持“打平了”，它只是没有变弱的迹象。要拿到结论得换更强的对手阵容，或者改成看“偏离引擎的程度”这类指标，而不是看胜负。

有一条顺序保证是刻意做成结构而不是提示词约束的：每回合开始时先让对手独立做决定——它读的是回合前的公开局面——然后才开放我方的操作面板。它决定的时刻早于我的选择，所以不可能参考我的选择。对方面板会直接写明“已独立出招（看不到你的选择）”，按钮是禁用的，不是点了没反应。

也顺手记一个由此产生的设计问题。“电脑对手”和“训练模式”在玩法上其实是同一个东西——同一套 AI、同一套结算，区别只在教练闭不闭嘴。这个区别是真的，但我原来把它藏在名字后面，第一次给用户看的时候被直接问了“这两个有啥区别”。现在两个入口的名字直接写明差别：训练是“按关卡挑战”，对局是“对手自由配队”。

这一轮抓到两个真 bug，两个都不是单元测试发现的。

第一个：训练模式下 \code{decideEnemyFirst} 只在分屏那条分支被重新调用，于是从第 2 回合起复用了上一回合锁定的行动。表现是对手一直出同一手，对局打不完。单测全绿，因为被复用的那个行动本身完全合法。它和本节开头那六个是同一个形状——测试测的是函数对不对，不测函数有没有在正确的时机被调用。

第二个更该记：玩家补位的那一回合本来就不该出招，却\textbf{被统计成了一次“引擎兜底”}。代码没坏，是口径错了，而错的方向恰好让这套 agent 显得比实际更差。发现它之后我把补位回合从统计里摘掉，前面“兜底 0 次”才是干净的。测量本身也会说谎，这一条比第一个 bug 更值得记。

\subsection{教练在 PVP 里该不该说话，我改错了两次}

这一条是我这次最值得记的。

题目对军师的描述是“在对战 / 阵容场景给出决策建议”，而场景设定写的是 PVP 对战。也就是说军师本来就该在对战里工作。

我第一版的做法是：只要是对局模式，教练一律不介入。理由是“两人共用一块屏，给一方提示就是不公平”。这个理由听起来成立，我还把“公平性设计”写进文档当亮点。问题是这等于把主功能在题目指定的场景里禁用了。改法是把判据从“是不是对局模式”换成“有没有第三方受影响”：对手是真人时闭麦，对手是电脑时照常工作。

改完又有人指出这个版本也不自洽：电脑对手存在的意义就是模拟真人，那它和真人局的行为就该一致，否则这个模拟本身就是假的。这句话是对的。最终改成本地对战一律允许教练，只保留线上竞技的闭麦开关。真正需要守住的底线不是“整局不说话”，而是“不读取对方待执行的动作”——那条项目一直没破，也有测试。

前后改了三次，前两次都是我把自己的推论当成了结论。

既然本地对战两边都能用教练，公平性就不需要靠禁用某个功能来实现了。做法是双方各有一条教练条，跑同一套引擎、同一套规则，各自只分析自己那一侧。对面用的是把 player 和 enemy 对调的镜像上下文，所以它看到的是对方的局面，也读不到任何隐藏信息——待执行的动作本来就不在上下文里。实测两条教练在同一回合给的结论不同，而且各自正确：一方建议“火花”，另一方建议“换上晶角犀”。

分屏这一块还踩过一个坑：两侧走不同的渲染函数，右边多一行标题，导致右侧的标签、按钮、教练条整体错位约 30 像素。改成两侧共用同一个渲染函数之后对齐了。

\subsection{换一张表要沿着所有抄写它的地方走一遍}

原来的属性相克表是不对称的：岩克三个属性，火和风各克一个，雷只被一个克。我数完之后发现说不出它为什么是这个形状——它不是设计出来的，是攒出来的。

换成两个独立的三环：火克草、草克水、水克火；岩克雷、雷克风、风克岩；环与环之间中性。现在每个属性恰好克一个、被一个克，这是一条可以直接断言的规则性质，不用靠眼睛看。同系伤害也从 $\times 0.75$ 改成 $\times 1$：原来它和“打向克制你的属性”共用了同一个分支，但这两件事不是一回事——游戏设定里没有任何东西支持“打自己人伤害更低”。

换完之后我以为最麻烦的是代码，结果代码是最省事的：军师、经验层、规则页都从同一张表读，跟着就走了。不跟着走的是那些把规则抄成自然语言的副本。手写的那张边界卡里写着旧规则，我改掉了；扫第二遍才发现同一类问题换了一层还在——六张由脚本生成的属性卡里仍然写着“同属性或被反克时 0.75”，而它们正是检索喂给模型的地面事实。也就是说教练会照着一张引擎已经不再实现的规则回答，而测试全绿，因为没有任何测试盯这句话。

根因是那句话在生成脚本里是硬编码的字符串，不在引擎的表里。这一轮把它改完了：\code{scripts/build-knowledge.js} 现在从 \code{RULES.typeAdvantage} 与 \code{RULES.typeResist} 取数，生成出来的六张卡直接写着“克制 1.5；打向克制你的属性 0.75；其余（含同系）1”。改引擎即自动同步，不再有第二份可以写错的副本。

值得留下的是这条教训的形状：改一个常量时真正要跟着改的，不是引用它的代码——代码会自动跟着走——是那些把规则抄成自然语言的副本，而它们没有类型系统盯着，也没有测试盯着。手工找这类副本我漏了三次，最后一次是知识卡自己的 principle 正文。这类问题不该靠人眼找，该靠生成器只从一处取数。

平衡矩阵也是这轮用新表重跑的，数字在下一节。
